\documentclass[11pt, letterpaper]{article}

\usepackage[utf8]{inputenc}
\usepackage[T1]{fontenc}
\usepackage[margin=0.85in]{geometry}
\usepackage{amsmath}
\usepackage{amssymb}
\usepackage{graphicx}
\usepackage{booktabs}
\usepackage{array}
\usepackage{float}
\usepackage{caption}
\usepackage{microtype}
\usepackage{enumitem}
\usepackage[numbers]{natbib}
\usepackage[colorlinks=true, linkcolor=blue, citecolor=blue, urlcolor=blue]{hyperref}

\captionsetup{labelfont=bf, font=small, skip=2pt}
\setlist{itemsep=2pt, parsep=2pt, topsep=3pt}
\setlength{\textfloatsep}{6pt}
\setlength{\intextsep}{6pt}
\setlength{\floatsep}{6pt}

\title{\textbf{Agentic Strategic Oversight of Classical Pathfinding with Cross-Episode Episodic Memory}\\[0.4em]
       \large }
\author{Fatih Karatay\\[0.2em]
        \small EN.705.625.8VL.SP26 Introduction to Agentic AI }
\date{April, 2026}

\begin{document}

\maketitle
\vspace{-2em}

\begin{abstract}
Classical dynamic-replanning algorithms such as A* and D* Lite efficiently navigate discrete grids but exhibit two properties that limit their use in real-world autonomous-navigation deployments: they are stateless across episodes. Each run begins without any knowledge of prior traversals and they produce no human-readable rationale for the trajectories they choose. I propose and evaluate a hybrid agentic system in which a LangGraph orchestrator acts as a strategic supervisor over an A* planner, augmenting it with two capabilities classical replanning cannot provide: (1) a persistent hazard map updated after every episode by a Reflection Node, encoding environment-specific risk patterns over time, and (2) an LLM (gpt-4o-mini via the OpenAI API) that produces natural-language risk assessments and end-of-episode reflections, creating a complete audit trail. I compare four system variants: Pure A*, D* Lite, an Agent-NoReflection ablation, and the Full Agent across three environment difficulty tiers (Static, Dynamic-Low, Dynamic-High) on a 20×20 grid over 50 sequential episodes per condition (600 episodes total). D* Lite achieves the highest success rate on Dynamic-Low (48\%) versus the Full Agent (38\%); reflection adds a consistent 2 pp over the no-reflection variant. The agent is approximately 1,500× slower than D* Lite per episode but produces interpretable justifications for every decision. I frame this as an honest cost-quality tradeoff: when raw success matters, classical replanning wins; when interpretability matters, the agent delivers an artifact no algorithmic baseline produces.
\end{abstract}

\section{Introduction}

Pathfinding in dynamic, partially-observed environments is a foundational problem in autonomous navigation, robotics, and game AI. Hart, Nilsson, and Raphael's A* algorithm [1] guarantees optimal-length paths in static graphs, and Koenig and Likhachev's D* Lite [2] extends this to incremental replanning when edge costs change mid-execution. Both, however, share two properties that become limitations as soon as a deployed system needs to explain itself or learn from experience:

\begin{itemize}
    \item \textbf{Stateless across episodes.} D* Lite reuses search state within an episode but resets at episode boundaries; the same congested corridor that caused a near-miss yesterday is rediscovered from scratch tomorrow.
    \item \textbf{No interpretable output.} A planner returns a path, not a rationale. When a robot avoids a corridor, an operator cannot ask why only whether the choice was optimal under the cost function.
\end{itemize}

These limitations matter as soon as the deployed system has any human in the loop or any cross-time learning requirement. Recent work on LLM-augmented planners LLM+P [3] generates PDDL plans for classical solvers, SayPlan [4] uses scene graphs for high-level task planning has demonstrated that LLMs can act as a "strategic layer" over algorithmic engines. To my knowledge, no published system applies this pattern to trajectory-level risk assessment with cross-episode memory.

This project contributes:

\begin{enumerate}
    \item A hybrid LangGraph-based agent that supervises an A* planner with a persistent, cross-episode hazard map and an LLM-driven reflective reasoner.
    \item A four-variant ablation study (A*, D* Lite, Agent-NoReflection, Full Agent) on three environment difficulty tiers, with all variants run on identical seeds for direct comparison.
    \item An empirical characterisation of the cost--quality tradeoff this design implies, with the Full Agent's natural-language justifications and reflections retained as a primary deliverable.
\end{enumerate}

The rest of the paper is organised as follows. Section 3 situates the system within prior work and the PEAS framing. Section 4 describes the seven-node LangGraph architecture, the persistent hazard-map mechanism, and the experimental harness. Section 5 details the protocol; Section 6 presents results and analysis. Section 7 concludes.

\section{Background and Related Work}

\subsection{Theoretical grounding}
The agent's operating context can be expressed in PEAS terms. The performance measure combines path success rate, path-efficiency ratio (path length relative to a Manhattan lower bound), collision rate, replanning count, decision latency, and LLM token cost. The environment is a discrete N×N grid containing static walls, dynamic obstacles that move under a configurable motion model, and a probabilistic hazard map derived from cross-episode memory. The agent's actuators are path-override commands, replanning triggers, memory writes, and cost-map adjustments. Its sensors include grid-state snapshots, the current planned path, the hazard map, an episode log, and retrieved past-episode summaries.

The system is a hybrid deliberative-reactive agent: pure-function nodes (Environment Scanner, Execution Monitor) handle reactive sensing and stepping, while deliberative nodes (A* Planner, Path Healer, Reflection) make planning and learning decisions. The Path Risk Evaluator and Reflection Node are deliberative-LLM nodes. They delegate the synthesis to an external language model and parse structured JSON output.

\subsection{Related work}
\begin{table}[H]
    \centering
    \small
    \caption{Comparison of prior planning systems and the proposed approach.}
    \begin{tabular}{p{0.16\textwidth} p{0.30\textwidth} p{0.42\textwidth}}
        \toprule
        \textbf{System} & \textbf{Approach} & \textbf{How this work differs} \\
        \midrule
        A* [1] & Optimal search in static graphs & No replanning; no cross-episode learning \\
        D* Lite [2] & Incremental replanning in dynamic graphs & Stateless across episodes; no interpretable reasoning \\
        LLM+P [3] & LLM generates PDDL plans for classical solvers & Task-level only; no trajectory-level risk assessment \\
        SayPlan [4] & LLM + scene graph for 3D task planning & High-level spatial reasoning; no episodic memory \\
        This work & LLM strategic oversight + episodic memory + reflective reasoning over A* & Trajectory-level risk + cross-episode memory + audit trail \\
        \bottomrule
    \end{tabular}
\end{table}
The proposed system occupies a niche distinct from prior LLM-planning work: it targets trajectory-level risk assessment (not task-level decomposition), operates in episodic sequential settings (not single-run), and treats natural-language decision justification as a primary output rather than a side effect.

\section{System Design and Methods}

\subsection{Architecture}
The system is implemented as a directed graph of seven LangGraph nodes (Figure 1). All nodes share an AgentState TypedDict; each node returns a partial state update that LangGraph merges into the running state.

\begin{figure}[H]
    \centering
    \includegraphics[width=0.72\textwidth]{fig1_langgraph_topology.png}
    \caption{LangGraph agent topology. Solid edges are unconditional; dashed edges are routed by AgentState fields such as risk score (\texttt{risk\_score}), obstacle detected (\texttt{obstacle\_detected}), stuck count (\texttt{stuck\_count}), and episode status (\texttt{episode\_status}). Node colours encode the kind of reasoning each performs: reactive (blue), deliberative (orange), deliberative-LLM (purple), memory (green), and reflective (red).}
    \label{fig:langgraph-architecture}
\end{figure}

The seven nodes are summarised in Table~\ref{tab:seven-nodes}.

\begin{table}[H]
    \centering
    \small
    \caption{Roles of the seven LangGraph nodes in the proposed architecture.}
    \label{tab:seven-nodes}
    \begin{tabular}{p{0.20\textwidth} p{0.15\textwidth} p{0.57\textwidth}}
        \toprule
        \textbf{Node} & \textbf{Type} & \textbf{Role} \\
        \midrule
        Environment Scanner & Reactive & Reads the current grid; computes a per-step uncertainty overlay by blending real-time dynamic-obstacle proximity with the persistent hazard map. \\
        A* Planner & Deliberative & Wraps a custom A* implementation with the per-step uncertainty overlay as soft cost weights. Calls \texttt{plan()} on first invocation and \texttt{replan()} thereafter. \\
        Memory Manager (read) & Memory & Retrieves up to five most recent past-episode summaries from a JSON-Lines store, filtered by grid size and obstacle density. \\
        Path Risk Evaluator & Deliberative-LLM & Asks the LLM to score the planned path's risk and justify the score. Output: \texttt{\{risk\_score: float, justification: str\}}. Fail-open: defaults to 0.0 on any LLM failure. \\
        Path Healer & Deliberative & When risk exceeds the configured threshold, raises the hazard score on flagged waypoints with a bump-and-replan strategy. Loops control back to the planner for replanning. \\
        Execution Monitor & Reactive & Advances the agent one step along the planned path. Detects newly blocked waypoints and routes back to risk evaluation when needed. Tolerates short stuck periods. \\
        Reflection Node & Reflective-LLM & Runs at episode end. Asks the LLM to diagnose the episode and propose hazard cells worth remembering. Updates the persistent hazard map via the Memory Manager. \\
        \bottomrule
    \end{tabular}
\end{table}

The conditional routing implements the proposal's reasoning policy:

\begin{itemize}
    \item After Path Risk Evaluator: heal if \texttt{risk\_score} $\geq \theta$ and \texttt{heal\_count < max\_heals}; otherwise, execute.
    \item After Path Healer: always loop back to the planner for a replan with the modified cost map.
    \item After Execution Monitor: terminal status $\rightarrow$ Reflection; running with an obstacle blocked $\rightarrow$ Risk Evaluator (mid-episode oversight); stuck (empty path, world has stepped) $\rightarrow$ Planner (try replanning); otherwise $\rightarrow$ Execution Monitor (advance the next step).
\end{itemize}

\begin{figure}[H]
    \centering
    \includegraphics[width=0.85\textwidth]{fig2_data_flow.png}
    \caption{Per-step data flow. Each node receives the shared AgentState and returns a partial update with the fields it owns. The Path Healer and Reflection nodes are the only ones that mutate the LLM-influenced cost surface (\texttt{hazard\_map}) and the persistent memory (\texttt{grid.hazard\_map}), respectively.}
    \label{fig:data-flow}
\end{figure}

\subsection{Reasoning, Memory, and Learning}
Three mechanisms together implement the agent's adaptation:

\begin{enumerate}
    \item \textbf{Per-step uncertainty overlay (reactive).} The Environment Scanner blends \texttt{grid.hazard\_map} (the persistent memory) with real-time dynamic-obstacle proximity into a per-step \texttt{state hazard\_map} that the A* Planner uses as cost weights.
    \item \textbf{Bump-and-replan (deliberative).} When the Risk Evaluator flags a path as high risk, the Path Healer raises the hazard score on the flagged waypoints and loops back to the planner. A* replans within the modified cost space, preserving its optimality guarantees relative to the bumped weights.
    \item \textbf{Reflection writes to persistent memory (reflective).} At episode end, the Reflection Node asks the LLM to identify cells worth remembering, then writes them into \texttt{grid.hazard\_map} via \texttt{update\_hazard\_map(...)}. Because \texttt{grid.hazard\_map} survives \texttt{grid.reset()}, this constitutes the project's cross-episode memory mechanism --- the only learning signal the agent retains between episodes.
\end{enumerate}

\subsection{Uncertainty and Bounded Rationality}
The LLM operates on a compressed text summary of the path and hazard map (approximately 400 input tokens) rather than on raw matrices, keeping per-call cost tractable. Risk assessments are gated behind a configurable threshold, and the Path Healer is further bounded by \texttt{max\_heals = 3} to prevent thrashing. JSON output is parsed with defensive validation, including clamping \texttt{risk\_score} to the interval $[0, 1]$ and filtering out-of-bounds hazard cells. Any LLM error, including authentication, network, or malformed-JSON failures, returns a fail-open fallback so the agent degrades gracefully to A* plus hazard-map behaviour rather than crashing.

\subsection{Implementation}
The full system is approximately 1,800 lines of Python organised under src/ (planners, environment, agent nodes, memory store, LLM client). Each LangGraph node lives in its own file, satisfying the rubric's modularity requirement: every reasoning element is identifiable in the codebase. The grid environment, both classical planners (A* and D* Lite implemented to match Koenig \& Likhachev's original 2002 specification including the \texttt{k\_m} accumulator and lazy heap deletion), and all seven agent nodes are covered by 83 pytest cases. The total LLM dependency is the OpenAI Chat Completions endpoint with JSON-mode response parsing.

\section{Experiments and Evaluation}

\subsection{Environment configurations}
Three environment tiers were evaluated on 20$\times$20 grids, as summarised in Table~\ref{tab:environment-tiers}.

\begin{table}[H]
    \centering
    \caption{Environment tiers used in the evaluation.}
    \label{tab:environment-tiers}
    \begin{tabular}{p{0.22\textwidth} p{0.15\textwidth} p{0.38\textwidth}}
        \toprule
        \textbf{Tier} & \textbf{Density} & \textbf{Obstacle motion} \\
        \midrule
        Static & 15\% & None \\
        Dynamic-Low & 15\% & Random walk, 1 cell every 5 steps \\
        Dynamic-High & 20\%* & Random walk, every step \\
        \bottomrule
    \end{tabular}
\end{table}
\textit{*Density was relaxed from the proposal's 25\% to 20\% because I did not implement the proposal's "periodic spawning"; without it, 25\% routinely produced no-path episodes that did not exercise the agent's adaptive logic.}

\begin{figure}[H]
    \centering
    \includegraphics[width=0.65\textwidth]{fig3_grid_world.png}
    \caption{A representative Dynamic-Low grid after 15 steps. Light pink cells encode the persistent hazard map (cross-episode memory), darker grey cells are static walls, orange cells are dynamic obstacles in their current positions, and the dashed blue line shows an illustrative diagonal path. The Execution Monitor advances along this path one cell per step, with grid.step() advancing the world after each move.}
    \label{fig:grid-world}
\end{figure}

\subsection{Variants}
Variants are summarised in Table~\ref{tab:variants}.

\begin{table}[H]
    \centering
    \caption{System variants evaluated in the experiments.}
    \label{tab:variants}
    \begin{tabular}{p{0.24\textwidth} p{0.56\textwidth}}
        \toprule
        \textbf{Variant} & \textbf{Components active} \\
        \midrule
        A* & A* + Execution Monitor (no replanning on obstacle detection) \\
        D* Lite & D* Lite incremental replanner (Koenig \& Likhachev, 2002) \\
        Agent-NoReflection & Full LangGraph agent, Reflection Node disabled (template summary, raw memory) \\
        Full Agent & Full LangGraph agent, all 7 nodes active \\
        \bottomrule
    \end{tabular}
\end{table}

\subsection{Protocol}
Each (variant $\times$ condition) cell ran 50 sequential episodes with seeds 1000--1049. Identical seed sequences were used across variants for fair comparison. Within a cell, episodes ran sequentially (not shuffled) so cross-episode memory could accumulate across runs in agent variants. Episodes terminated on success (goal reached), collision (agent enters a blocked cell, which can occur when a dynamic obstacle moves onto the agent's current position), or timeout (steps $\geq 3 \times$ Manhattan distance). The configurable risk threshold was set to 0.3 and \texttt{max\_heals = 3}.

\subsection{Metrics}
The experiment runner emits one JSON record per episode capturing outcome, steps taken, replan count, heal count, wall-clock latency, LLM call count, total tokens, and Path Efficiency Ratio (PER = steps / Manhattan). The aggregator (\texttt{experiments/aggregate\_results.py}) reduces these records into the per-cell summary in Section 6.

\section{Results and Analysis}
\subsection{Headline results}
The summary across all 600 episodes is shown in Table~\ref{tab:headline-results}.

\begin{table}[H]
    \centering
    \caption{Success rates by system variant and environment tier.}
    \label{tab:headline-results}
    \begin{tabular}{p{0.28\textwidth} p{0.14\textwidth} p{0.18\textwidth} p{0.18\textwidth}}
        \toprule
        \textbf{Variant} & \textbf{Static} & \textbf{Dynamic-Low} & \textbf{Dynamic-High} \\
        \midrule
        A* & 96\% & 20\% & 0\% \\
        D* Lite & 96\% & 48\% & 6\% \\
        Agent-NoReflection & 96\% & 36\% & 0\% \\
        Full Agent & 96\% & 38\% & 0\% \\
        \bottomrule
    \end{tabular}
\end{table}

\begin{figure}[H]
    \centering
    \includegraphics[width=0.72\textwidth]{fig4_success_rate_by_variant.png}
    \caption{Success rate by variant and condition. Static is dominated by A*-grade planning (no dynamics to react to); the Dynamic-Low tier is where variants meaningfully separate; Dynamic-High is too chaotic for any single-step reasoning to consistently succeed.}
    \label{fig:success-rate}
\end{figure}

\subsection{Static}
All four variants achieve identical 96\% success. The 4\% failure (2 of 50 episodes) corresponds to seeds whose random static-wall layout happens to disconnect the start cell from the goal; A* finds no path and the episode times out. Because all variants encounter the same seeds, the failure is shared. The fact that the agent's overhead does not improve on this is the expected result: when the world does not move, A*-grade planning is sufficient and the agent's reasoning is wasted work.

A second observation in this tier is more subtle. The Full Agent triggers an average of 2.58 heals and 2.86 replans per static episode, purely from reflection-baked hazards accumulated in earlier episodes. Agent-NoReflection, which writes only static-template summaries to memory, triggers essentially no heals (0.00) on the same seeds. The Reflection Node's hazard-cell suggestions (Section 4) are therefore the dominant driver of the persistent hazard map, not the raw observed-obstacle bake. This is a clean qualitative confirmation that the LLM reflection is doing real work, even if its impact on success rate is modest.

\subsection{Dynamic-Low}
This is the headline tier. D* Lite (48\%) outperforms the Full Agent (38\%), and both clearly beat unmodified A* (20\%). The Reflection Node adds a consistent 2 pp over Agent-NoReflection (38\% vs 36\%), with the Full Agent triggering an average of 2.80 heals per episode (vs 2.30 for the no-reflection variant). The reflection signal is real but small.

The Full Agent's persistent hazard map produces longer paths (PER 1.18 vs D* Lite's 1.03) because flagged waypoints are detoured around. The longer paths mean the agent spends more time exposed to obstacle motion, generating more collision opportunities, which creates a classic over-correction failure mode. Lowering the risk threshold from 0.5 (a 20-episode pilot) to 0.3 in this run made the agent more conservative and worse by ~7 pp in success rate, but with measurably more triggered heals (0.70 to 2.80 per episode).

\subsection{Dynamic-High}
At 20\% obstacle density and every-step motion, only D* Lite achieves any successes (6\%). All other variants fail completely. The Full Agent does manage to convert some collisions into healed-and-replanned attempts (an average of 2.88 heals and 3.04 replans per episode), but cannot consistently navigate to the goal under these conditions.

\subsection{Episode-over-episode Learning}

\begin{figure}[H]
    \centering
    \includegraphics[width=0.78\textwidth]{fig5_learning_curve.png}
    \caption{Rolling success rate (window = 5) across the 50 sequential dynamic-low episodes. The Full Agent (green) climbs from 0 to 80\% around episodes 10--15 as the persistent hazard map accumulates obstacle-frequency signal, then declines as later episodes draw harder seeds. D* Lite (blue) shows comparable peak performance but no learning trend; its replanning is lifetime-bounded but episode-resetting. Agent-NoReflection (orange) tracks the Full Agent closely, suggesting most of the cross-episode benefit comes from raw memory writes rather than LLM reflection synthesis.}
    \label{fig:learning-curve}
\end{figure}

The learning curve provides direct evidence for cross-episode memory: the Full Agent's rolling success rate rises from 0 to a peak of ~0.8 over the first 15 episodes, indicating the persistent hazard map is shaping subsequent A* plans. D* Lite reaches comparable peaks but oscillates without a clear trend, consistent with its episode-stateless design. The drop in the latter half of the curve reflects natural seed variance rather than learning regression; agents and baselines decline together.

\subsection{Cost-quality Tradeoff}

\begin{figure}[H]
    \centering
    \includegraphics[width=0.78\textwidth]{fig6_cost_quality_scatter.png}
    \caption{Mean wall-clock per episode (log scale) versus success rate. The agent variants cluster in the 1,500--19,000 ms range, which is three to four orders of magnitude slower than D* Lite, without producing better success rates on any condition.}
\end{figure}

The Full Agent consumes a mean of 7,340 LLM tokens per dynamic-low episode and takes 18.9 seconds wall-clock, compared to D* Lite's 13 ms. This is approximately a 1,500× slowdown for a worse outcome on raw success rate. The agent's per-decision LLM justifications and end-of-episode reflections are the value being purchased: every routing decision the agent makes is recorded with a human-readable rationale, producing a complete audit trail per episode.

\subsection{Answering the Research Questions}

\begin{itemize}
    \item \textbf{RQ1 (Memory):} Cross-episode memory accumulation is empirically observable in the learning curve (Figure 5). The Full Agent climbs from 0 to approximately 80\% rolling success rate over the first 15 episodes. But does not in my setup translate into a higher aggregate success rate than D* Lite. The bump-and-replan strategy over-corrects on the relatively mild dynamic-low tier.
    \item \textbf{RQ2 (Reflection):} The reflection ablation produces a small but consistent 2 pp gain (38\% vs 36\%) on dynamic-low. The signal is real but modest at this scale; the proposal's anticipated 50-episode horizon is at the edge of what is needed to detect it cleanly.
    \item \textbf{RQ3 (Cost-Quality):} In my experiments the LLM oversight does not justify its latency or token cost on the basis of raw success rate. Its value must be measured on a different axis.
\end{itemize}

The proposal explicitly anticipated this outcome: "D* Lite may outperform the Full Agent on per-episode metrics in Dynamic-Low conditions; this is an expected and honest result that the paper will frame as a cost-quality trade-off rather than a failure." That framing holds.

\section{Conclusion and Future Work}
I built and evaluated a hybrid LangGraph-based agent that supervises an A* planner with cross-episode episodic memory and an LLM-driven risk and reflection layer. Across 600 episodes spanning four variants and three environment tiers, the Full Agent did not exceed D* Lite on raw success rate, and its mean wall-clock per episode is roughly 1,500× higher. However, the agent produces a complete audit trail --- natural-language risk justifications and end-of-episode reflections --- that no algorithmic baseline can offer, and demonstrably accumulates cross-episode hazard knowledge that influences future planning (Figure 5).

\subsection{What was Learned}
\begin{itemize}
    \item \textbf{Cross-episode memory works mechanically.} The Reflection Node successfully writes hazard cells into the persistent map, the rolling-success learning curve shows the expected accumulation pattern, and the Risk Evaluator's heal triggers track the accumulated signal.
    \item \textbf{Bump-and-replan over-corrects in mild-dynamics environments.} Each heal extends the path; longer paths increase exposure to obstacle motion. Lowering the risk threshold made the agent strictly worse.
    \item \textbf{LLM oversight is expensive.} Even at \texttt{gpt-4o-mini} pricing, the agent costs approximately \$1 USD across 600 episodes --- affordable, but the per-episode latency (approximately 18 seconds for a Full Agent on dynamic-low) is incompatible with real-time deployment without aggressive caching or model distillation.
    \item \textbf{Interpretability is the agent's defensible contribution.} Every routing decision is captured with a rationale; every episode ends with a synthesised diagnosis stored to disk.
\end{itemize}


\begin{thebibliography}{9}

\bibitem{hart1968}
P. E. Hart, N. J. Nilsson, and B. Raphael.
\newblock A formal basis for the heuristic determination of minimum cost paths.
\newblock \textit{IEEE Transactions on Systems Science and Cybernetics}, 4(2):100--107, 1968.

\bibitem{koenig2002}
S. Koenig and M. Likhachev.
\newblock D* Lite.
\newblock In \textit{Proceedings of the AAAI Conference on Artificial Intelligence}, 2002.

\bibitem{liu2023}
B. Liu et al.
\newblock LLM+P: Empowering large language models with optimal planning proficiency.
\newblock \textit{arXiv preprint} arXiv:2304.11477, 2023.

\bibitem{rana2023}
K. Rana et al.
\newblock SayPlan: Grounding large language models using 3D scene graphs for scalable robot task planning.
\newblock In \textit{Conference on Robot Learning (CoRL)}, 2023.

\bibitem{langgraph}
Anthropic / LangChain.
\newblock LangGraph: State-machine workflows for LLM agents.
\newblock \url{https://github.com/langchain-ai/langgraph}

\bibitem{openai-chat}
OpenAI.
\newblock Chat Completions API and gpt-4o-mini model.
\newblock \url{https://platform.openai.com/docs/api-reference/chat}

\end{thebibliography}



\end{document}